\fichatitulo{58 · GOVERNANÇA}{Government Information Quarterly · 2025}{Governança de IA em organizações públicas}
{\small DE ALMEIDA, Patricia Gomes Rêgo; SANTOS JÚNIOR, Carlos Denner dos. \textit{Artificial intelligence governance: Understanding how public organizations implement it}. Government Information Quarterly, v. 42, 102003, 2025. \href{https://doi.org/10.1016/j.giq.2024.102003}{DOI} e PDF consultado.\par}

\campo{Problema e objetivo}
Como organizações públicas incorporam princípios éticos, políticas e práticas de governança de IA em seus processos e sistemas.

\campo{Principal contribuição · síntese interpretativa}
Propõe o framework AIGov4Gov a partir da análise de práticas de governança em organizações públicas. O estudo relaciona governança a níveis estratégico, tático e operacional e destaca treinamento de decisores, desenvolvedores e usuários.

\campo{Método / natureza da evidência}
Estudo exploratório e descritivo de 28 organizações públicas em cinco continentes, combinando questionários, entrevistas, análise documental e QCA crisp-set e fuzzy-set.

\campo{Resultado principal}
Organizações que treinam decisores, desenvolvedores e usuários apresentam combinações associadas a estágios mais avançados de governança. O artigo também identifica estratégia, políticas, princípios éticos, padrões e processos de auditoria como componentes do framework proposto.

\campo{Excertos-chave · seleção literal}
\begin{itemize}
\excerto{Amostra}{28 public organizations across five continents}{Resumo, p. 1.}
\excerto{Treinamento}{training decision-makers, AI system developers, and users}{Resumo, p. 1.}
\end{itemize}

\campo{Limitações e alcance}
\textbf{Fonte:} o estudo é exploratório, a amostra não reproduz a distribuição global de organizações e a análise não aprofunda a qualidade de cada treinamento. \textbf{Análise da ficha:} sustenta governança e capacitação institucional, não desempenho de RAG nem causalidade geral.

\campo{Aproveitamento na dissertação · proposta}
\textbf{Destino:} referencial teórico e discussão. \textbf{Uso:} justificar supervisão, capacitação e processos de governança para interpretar a avaliação de IA em organização pública.

\registro{Fonte consultada: PDF local, 18 páginas; resumo, introdução, método, resultados de QCA, framework AIGov4Gov e limitações. Chave: \texttt{deAlmeidaSantos2025aiGovernance}.}
